% ---
% id: receiver/rem:n0
% kind: rem
% title: "The monotonicity step at zero photons"
% summary: >-
%   The monotonicity step of [[lem:collapse]] fails at n = 0, with an explicit
%   two-port receiver as witness. Nothing is lost, because the infimum runs
%   over photon numbers above the cutoff, all at least one.
% prominence: standard
% status: accepted
% uses: [lem:collapse]
% assumes: []
% refs: []
% provenance: newly-derived
% lean:
% ---

\begin{remark}
\label{rem:n0}
Lemma~\ref{lem:collapse} is stated from $n\ge1$ and cannot be extended to
$n=0$. Take $K\ge2$, $\gamma=\beta_k=1$ and $d_k=1$ for every $k$, and $c<1$: then
$\Bnd(0)=1+(K-1)-K=0\ge0$, while
$\Bnd(1)=1+(K-1)c-K=(K-1)(c-1)<0$, so \eqref{eq:collapsestep} fails at $n=0$.
The data are realised, for instance, by $K=2$ on two input modes with
$A_1=a\dyad{0}$, $A_2=a\dyad{1}$, $a\in(0,1)$: there $c=1-a$ and
$d_1=d_2=1$. The substitution $u_k=\beta_kd_k^{\,n}$ degenerates at $n=0$,
which is precisely where the proof breaks. Nothing is lost: the infimum in
\eqref{eq:collapseinf} runs over $n>N\ge0$, so the induction never starts below
$n=1$.
\end{remark}
